% Generated draft: requires mathematical and source review.
\documentclass[11pt]{article}
\usepackage[margin=25mm]{geometry}
\usepackage{fontspec}
\setmainfont{Latin Modern Roman}
\usepackage{amsmath,amssymb,mathtools,unicode-math}
\setmathfont{Latin Modern Math}
\usepackage{enumitem,needspace,xurl,hyperref}
\setlength{\parindent}{0pt}
\setlength{\parskip}{5pt}
\setlength{\emergencystretch}{3em}
\newcommand{\N}{\mathbb{N}}
\newcommand{\Z}{\mathbb{Z}}
\newcommand{\Q}{\mathbb{Q}}
\newcommand{\R}{\mathbb{R}}
\newcommand{\C}{\mathbb{C}}
\newcommand{\F}{\mathbb{F}}
\newcommand{\A}{\mathbb{A}}
\newcommand{\PP}{\mathbb{P}}
\newcommand{\E}{\mathbb{E}}
\newcommand{\eps}{\varepsilon}
\newcommand{\dd}{\,\mathrm{d}}
\begin{document}
\section*{UnsolvedMath ID: 30005459}
\textbf{Source problem code:} catalogue entry
\subsection{Definitions and mathematical statement}
The supplied record concerns positivity of noncommutative state polynomials on a set denoted by \(K_C\), assumed convex in the source. It does not define the constraint data \(C\), the construction of \(K_C\), the algebra of state polynomials, or their evaluations. Consequently, the following is a conditional formulation of the two questions; these missing definitions must be recovered before it becomes a complete mathematical statement.

For bookkeeping only, let \(\mathcal P\) denote the source's intended class of noncommutative state polynomials, let \(\mathcal E(K_C)\) denote its intended admissible evaluations on \(K_C\), and let \(\mathcal M(K_C)\subseteq\mathcal E(K_C)\) denote its admissible matrix evaluations. Write \(\operatorname{Pos}(p,e)\) for the source's intended assertion that the evaluation of \(p\) at \(e\) is positive. These symbols do not specify the missing algebra, admissibility conditions, or positivity convention.

\begin{enumerate}
\item For the intended convex \(K_C\), are matrix evaluations sufficient to test positivity? The proposed sufficiency has the schematic form
\[
\forall p\in\mathcal P,\qquad
\left[\forall e\in\mathcal M(K_C),\ \operatorname{Pos}(p,e)\right]
\Longrightarrow
\left[\forall e\in\mathcal E(K_C),\ \operatorname{Pos}(p,e)\right].
\]
The reverse implication follows from the stipulated inclusion of evaluation classes. The source excerpt does not specify the allowed matrix sizes, the underlying scalar field, the state data, or whether positivity means nonnegativity or strict positivity. It also does not specify the range of constraint systems \(C\) covered by the question.
\item Is there an analogue of the Krivine--Stengle Positivstellensatz for \(K_C\) when evaluations on unbounded self-adjoint operators are considered? The supplied excerpt specifies neither the admissible operator tuples and states nor the common domains needed to evaluate polynomials, the meaning of positivity for those evaluations, or the algebraic certificate sought. No particular certificate identity or equivalence can therefore be stated from this record alone.
\end{enumerate}

The source motivates these questions by saying that its Theorem~1 characterizes positive state polynomials but does not characterize noncommutative state polynomials. The distinction between those classes and the hypotheses of that theorem are not supplied. The excerpt also says that its Theorem~3 yields an efficient state polynomial optimization procedure despite the two open questions; that theorem and the meaning of efficiency are not supplied.
\subsection{Short English statement}
The questions ask whether matrix evaluations detect positivity of noncommutative state polynomials on the source's convex set \(K_C\), and whether a Krivine--Stengle-type algebraic positivity characterization exists when unbounded self-adjoint operator evaluations are allowed. Essential definitions and operator-domain conditions are missing from the supplied excerpt, so this statement-only draft requires source consultation and human mathematical review. The supplied record labels both questions open; that status has not been independently verified.
\subsection{Sources}
\begin{itemize}
\item Supplied local UnsolvedMath registry record, permanent ID \(30005459\), title: \emph{Positivity of Noncommutative State Polynomials on Convex Sets}; source problem code \texttt{catalogue entry}; collection: \emph{Oberwolfach Reports Open Problems}. The record gives proposed year 2023 and identifies no individual proposer.
\item \emph{Real Algebraic Geometry with a View toward Koopman Operator Methods} (2023), Oberwolfach Reports, pp.~741--811, DOI 10.4171/owr/2023/14. The supplied extracted passage identifies p.~773. Source URL and supplied literature-source link, labelled ``original Oberwolfach report'': \url{https://doi.org/10.4171/owr/2023/14}.
\item Literature assessment dated \texttt{2026-08-22}: the supplied curation labels the questions open and reports that no primary result located in that review resolves either question in the stated noncommutative state-polynomial setting. This is supplied dated curation, not independently verified literature assessment. The background expressly requests independent verification of time-sensitive status and cited claims.
\item The registry assigns statement status \texttt{corrected\_verified} and reports that the official report was checked directly and that definitions were restored in the clean statement. These are supplied provenance assertions, not verification performed for this draft; the supplied clean statement still does not contain the definitions needed for a complete formulation.
\end{itemize}
\end{document}
